\documentclass[12pt]{article}
\usepackage[utf8]{inputenc}
\usepackage[T1]{fontenc}
\usepackage{hyperref}
\title{Sedoa Language Resources for K-12 Curriculum Development}
\author{Compiled by Claude Sonnet 5 (claude-sonnet-5) for Andr\'as Kornai}
\date{}
\begin{document}
\maketitle

\section{Introduction}
Sedoa (ISO 639-3 tvw, Glottocode sedo1237) is an Austronesian language of Central Sulawesi, Indonesia. It is usually classed in the Kaili--Pamona branch of Celebic. Wikipedia gives about 900 native speakers (2009) of an ethnic group of roughly 3,000. Sedoa is not one of the Badaic languages (Bada, Besoa, Napu), although they are spoken in the same region. Resources for those languages must not be mixed up with Sedoa.

\section{Primary Educational Resources}
None were found in search results.

\section{Linguistic and Reference Materials}
\begin{itemize}
\item Wikipedia, Sedoa language: \url{https://en.wikipedia.org/wiki/Sedoa_language}
\end{itemize}
No dictionary, grammar or orthography was found.

\section{Religious and Community Resources}
\begin{itemize}
\item Joshua Project, Sedoa (Tawaelia) in Indonesia, a weak source for classification: \url{https://joshuaproject.net/people_groups/14812/ID}
\end{itemize}

\section{Digital and Audio Resources}
None were found in search results.

\section{Curriculum Development Context}
The language is not known to be taught in schools. The searches found almost nothing beyond encyclopedic and people-group entries. Any curriculum would depend on new documentation work with the community, and on intergenerational transmission, since not all young people speak the language.
\end{document}